% Generated draft: requires mathematical and source review.
\documentclass[11pt]{article}
\usepackage[margin=25mm]{geometry}
\usepackage{fontspec}
\setmainfont{Latin Modern Roman}
\usepackage{amsmath,amssymb,mathtools,unicode-math}
\setmathfont{Latin Modern Math}
\usepackage{enumitem,needspace,xurl,hyperref}
\setlength{\parindent}{0pt}
\setlength{\parskip}{5pt}
\setlength{\emergencystretch}{3em}
\newcommand{\N}{\mathbb{N}}
\newcommand{\Z}{\mathbb{Z}}
\newcommand{\Q}{\mathbb{Q}}
\newcommand{\R}{\mathbb{R}}
\newcommand{\C}{\mathbb{C}}
\newcommand{\F}{\mathbb{F}}
\newcommand{\A}{\mathbb{A}}
\newcommand{\PP}{\mathbb{P}}
\newcommand{\E}{\mathbb{E}}
\newcommand{\eps}{\varepsilon}
\newcommand{\dd}{\,\mathrm{d}}
\begin{document}
\section*{UnsolvedMath ID: 10800014}
\textbf{Source problem code:} catalogue entry
\subsection{Definitions and mathematical statement}
Fix a positive odd integer \(d\). Let \(V_d=\R[x,y,z]_d\) be the real vector space of homogeneous degree-\(d\) polynomials and \(\PP(V_d)\) its coefficient projective space, so scalar multiples define the same equation. In the usual complex-nonsingular convention, set
\[
 \mathcal U_d^{(1)}=\bigl\{[F]\in\PP(V_d):
 C_F(\C)\subset\C\PP^2\text{ is nonsingular, and }
 C_F(\R)\subset\R\PP^2\text{ has exactly one connected component}\bigr\},
\]
where \(C_F\) is the zero locus of \(F\), and give this set the topology induced by the real coefficient projective space. Two such curves are rigid isotopic if they can be joined by a continuous path in this nonsingular parameter locus.

The question is whether \(\mathcal U_d^{(1)}\) is connected for every positive odd \(d\), equivalently whether every two members of this fixed-degree class are rigid isotopic. Along a nonsingular real family the number of real components is constant, so the path can be required to stay in \(\mathcal U_d^{(1)}\). The excerpt does not say whether nonsingularity is required over \(\C\) or only at real points; if the latter is intended, the parameter locus and allowed rigid isotopies must instead be adjusted accordingly. Different odd degrees are separate questions, not one coefficient-space component.
\subsection{Short English statement}
For each fixed odd degree, are all nonsingular real projective plane curves with exactly one real component rigid isotopic?
\subsection{Sources}
\begin{itemize}
\item Victor Vassiliev, \emph{A Few Problems on Monodromy and Discriminants}, Problem 5D. \url{https://amj.math.stonybrook.edu/html-articles/Files-2015-2024/15-11/}. Recorded access 2026-07-29; extraction html-structure.
\item Local UnsolvedMath registry, \texttt{data/unsolvedmath/problems.json}, permanent ID \(10800014\), exact source code \texttt{catalogue entry}. Its partially\_solved/PARTIAL-PROGRESS summary reports complete low-degree rigid-isotopy classifications, one-component connectedness in known cases, odd degrees \(d\leq5\), and the even-degree \(d=6\) context of 64 sextics, without separate primary citations. It reports the general odd-degree problem as remaining open.
\end{itemize}
\end{document}
